\answer{ex:19:1}{(a)~Một khớp cho $6.7\mV$ ($R=66.7$~M$\Omega$), nên để tới ngưỡng nói chung cần ít nhất $4$ khớp (ba khớp chỉ đủ theo tiệm cận); $8$ khớp để tới trong $10\ms$; $14$ khớp để tới trong $5\ms$. (b)~$0.1\ms\ll\tau_m$, hiệu chỉnh do rò $\sim0.25\%$.}
\answer{ex:19:2}{(a)~$1.75\cdot10^{10}$ (một phần tư số bộ trộn đáp ứng với mọi thứ), $4.4\cdot10^9$ và $0.06$ (với $k=40$ hầu như không bao giờ đáp ứng). (b)~$10^3$ lần: $2\cdot10^5$ so với $200$.}
\answer{ex:19:3}{(a)~$C(2n,n)=2^{2n-1}$ theo tính đối xứng của hệ số nhị thức. (b)~$0.93$ và $0.09$; độ rộng vùng chuyển tiếp theo $P$ là $\sim\sqrt{2n}$, tương đối là $\sim1/\sqrt n$.}
\answer{ex:19:4}{Từ một đuôi gai, $V_{\text{thân}}<\kappa E_E<\theta$ với bất kỳ số đầu vào nào; với $g_0=0.5g_L$ cần $n\ge3$ trên mỗi đuôi gai.}
\answer{ex:19:5}{$I\ge C\sigma_V/\sigma_t=15$~nA, tức $150$ khớp đồng bộ, không thực tế; nơ-ron nhỏ chỉ cần $1$~nA, và với một khớp $4$~nA độ rung pha là $5$~\textmu s.}
\answer{ex:19:6}{Cực đại ở đầu gần: $0.03\,\tau_m$ và $0.97$ ($L=0.2$); $0.32\,\tau_m$ và $0.66$ ($L=1$); $0.79\,\tau_m$ và $0.33$ ($L=2$). Ước tính~\eqref{eq:dendriteload} theo điện dung toàn phần chỉ đúng khi $L\ll1$.}
\answer{ex:19:7}{Bài mở. Khi $m$ cố định, thời gian đáp ứng tăng tuyến tính theo $P$ được ghi nhớ; khi tỉ lệ $m=\varphi n$ cố định, nó không còn phụ thuộc $n$, và sự chậm chạp của nơ-ron lớn là hệ quả của việc cộng nhiều đầu vào chứ không phải của kích thước tự thân.}
